\section*{章末問題}
\begin{enumerate}[label=\textbf{\thechapter.\arabic*},leftmargin=*]

 \item\label{問:極限と漸化式} 次の極限を求めよ.また,\,$a_1=1$,\,$a_{n+1}=\sqrt{2+a_n}$ で定まる数列が単調増加かつ上に有界であることを示し,その極限を求めよ.
 \[
  \text{(1)}\quad\lim_{n\to\infty}\frac{2n^2+3}{5n^2-n},
  \qquad
  \text{(2)}\quad\lim_{n\to\infty}\bigl(\sqrt{n^2+2n}-n\bigr).
 \]
 \item\label{問:等比誤差} $\sum_{k=0}^{\infty}\left(-\frac12\right)^k$ の和を求めよ.\,$k=0$ から $n-1$ までの和で近似するとき,誤差の絶対値を $10^{-3}$ 未満にする最小の自然数 $n$ を求めよ.
 \item\label{問:連続と微分可能} $x\ne0$ で $f(x)=\frac{\sin2x}{x}$,\,$x=0$ で $f(0)=c$ とする.\,$f$ が $0$ で連続となる $c$ を求めよ.また,\,$g(x)=|x-1|$ は $1$ で連続だが微分可能でないことを,差商を使って説明せよ.
 \item\label{問:微分計算} 次の関数を微分せよ.(2) については $x=0$ における接線も求めよ.
 \[
  \text{(1)}\quad\log(1+x^2),\qquad
  \text{(2)}\quad(x^2+1)\e^{-x},\qquad
  \text{(3)}\quad x^{\sin x}\quad(x>0).
 \]\end{enumerate}
\smallskip\noindent 解答は\pageref{解答:第2章}ページ.
